\section{Chapter~\ref{sec:sensors}. The Machine's Inputs}

The structure of the sensors has been measured directly, by recording currents from single cells and vibrations of the cochlea and by counting retinal cells. The sensitivity limit and the shot-noise formula follow from physics, while the claim that the sensors' codes are tuned for the greatest information is a hypothesis with strong support.

{\footnotesize
\setlength{\tabcolsep}{3pt}\renewcommand{\arraystretch}{1.15}
\begin{longtable}{>{\raggedright}p{0.5cm}>{\raggedright}p{4.0cm}>{\raggedright}p{5.6cm}>{\raggedright}p{2.3cm}>{\raggedright\arraybackslash}p{1.7cm}}
\toprule
No. & Claim & Experiment and what was measured & Source & Status \\
\midrule\endfirsthead
\toprule
No. & Claim & Experiment and what was measured & Source & Status \\
\midrule\endhead
1 & A sensor is a transducer: the receptor gives a current, and the output of the sensory cells of the eye and ear is glutamate onto the \nt{GLUT} bus; the first cells produce no spikes (Fig.~\ref{fig:transducer}) & Turtle rods and cones release an excitatory amino acid, detected by glutamate channels on an excised membrane patch. Inner hair cells of the rat cochlea: currents in the nerve-fiber endings flow through glutamate AMPA receptors, and the release rate grows with smooth depolarization of the cell up to $150$\,Hz & \cite{CopenhagenJahr1989,GlowatzkiFuchs2002} & measured \\
2 & A light sensor in the dark responds to a single photon with a current of about a picoampere & Suction electrode on the outer segment of a toad rod: responses to dim flashes are quantized, an event of $\approx1$\,pA ($5\%$ of the dark current) with a spread of $0.2$\,pA, efficiency $0.5$. Humans report the arrival of a single photon more often than chance & \cite{Baylor1979,Tinsley2016} & measured \\
3 & A sound sensor detects a displacement of less than a nanometer & M\"ossbauer technique, basilar membrane of the guinea-pig cochlea at the $16$--$19$\,kHz place: at the auditory-nerve response threshold the membrane velocity is about $0.04$\,mm/s, i.e., a displacement of $\approx0.35$\,nm (our conversion). The membrane was measured, not the sensor's bundle & \cite{Sellick1982,Hudspeth1989} & measured \\
4 & Precision in the dark is limited by photon shot noise~\eqref{eq:photonnoise}; in dim light integration is longer & The formula follows from Poisson statistics. In a rod the noise in dim light matches the sum of independent quantal events; against a bright background the response to a flash is smaller and shorter. The ``tenths of a second'' figure is not checked separately here & \cite{Baylor1979} & theory \\
5 & The input range is ten orders of magnitude for light and twelve for sound, while the spike rate spans less than three & For sound: the basilar-membrane response at the characteristic frequency grows with a slope as low as $0.2$\,dB/dB over $40$--$80$\,dB, and compression is seen even above $100$\,dB. The order-of-magnitude figures themselves are standard estimates, unsourced in the chapter & \cite{Ruggero1997} & measured \\
6 & The sensor response is~\eqref{eq:nakarushton}, and $\sigma$ follows the mean brightness, shifting the curve along the logarithmic axis (Fig.~\ref{fig:adaptation}) & The formula was first fitted to S-potentials of the fish retina. Turtle cones: responses obey $V=I V_m/(I+\sigma)$ and span $\sim3.5$ orders of magnitude of brightness; after $2$--$3$ minutes of background the curve shifts horizontally, keeping its width. The link to division by total activity is from a review & \cite{NakaRushton1966,NormannPerlman1979,Carandini2012} & measured \\
7 & Sensors transmit changes, and their codes are tuned for the greatest information per spike (Proposition~\ref{prop:sensors}) & The principle was proposed theoretically. The strongest evidence: in the fly, the ``contrast $\to$ response'' curve of first-layer neurons matches the cumulative distribution of contrasts in natural scenes, which yields the greatest information & \cite{Barlow1961,Laughlin1981} & hypothesis \\
8 & On and off sensors form a two-wire code; the event camera reproduces it in silicon & A review of experiments: the on and off channels of the retina separate already at the first synapse and can be switched off separately. Camera: $128\times128$ pixels without frames, $120$\,dB range, $15$\,\textmu s latency & \cite{Schiller1992,Lichtsteiner2008,Gallego2022} & measured \\
9 & The eye has $\sim10^8$ light sensors and $\sim10^6$ output neurons: compression by $\sim100$ & Counts on whole-mount human retinas: on average $92$ million rods and $4.6$ million cones; $0.7$ to $1.5$ million output (ganglion) cells in different eyes & \cite{Curcio1990,CurcioAllen1990} & measured \\
10 & The mouse eye has more than thirty output channels & Mouse: two-photon calcium imaging of more than $11\,000$ output cells and clustering of the responses gives clearly more than $30$ functional types. The number for humans has not been measured & \cite{Baden2016} & measured \\
11 & The guinea-pig eye transmits $\sim875\,000$ bits/s; scaling to the human eye gives $\sim10^7$ bits/s & Guinea pig, multielectrode array, motion in natural scenes: $6$--$13$ bits/s per cell, $\sim875\,000$ bits/s for $\sim10^5$ output cells. For humans ($\sim10^6$ cells), $\sim10^7$ bits/s is the authors' extrapolation & \cite{Koch2006} & measured \\
12 & The ear is a bank of bandpass filters with a nearly logarithmic frequency map (Fig.~\ref{fig:filterbank}) & The place of maximum basilar-membrane vibration depends on frequency; the frequency--place function is nearly exponential and a single form describes the cochleas of humans and living animals & \cite{Greenwood1990,Robles2001} & measured \\
13 & The resonators are active: some sensors amplify weak vibrations thousands of times in amplitude ($66$--$76$\,dB), and the gain falls as loudness grows & Laser vibrometry at the base of the chinchilla cochlea: for weak tones at the characteristic frequency the gain is $66$--$76$\,dB relative to the stapes; death lowers sensitivity by $60$--$81$\,dB (a factor of $10^3$--$10^4$ in amplitude) and removes the compression. The source of the force is the outer hair cells & \cite{Ruggero1997,Robles2001} & measured \\
14 & The ear's amplifier sits at the edge of self-oscillation & Calculation: near a Hopf bifurcation point, range compression, sharp tuning, and combination tones, all the known nonlinearities of hearing, are unavoidable. That the cochlea is tuned in exactly this way has not been proved directly & \cite{Eguiluz2000} & hypothesis \\
15 & At low frequencies spikes are phase-locked to the sound (in the guinea pig, locking weakens from $\sim600$\,Hz and is noticeable up to $\sim3.5$\,kHz), and direction is found from them; at high frequencies only the place code remains & Guinea-pig auditory nerve: phase locking begins to fall at about $600$\,Hz and disappears above $3.5$\,kHz; in the cat and the monkey the limit is about an octave higher. The interaural time difference is from a review & \cite{PalmerRussell1986,Grothe2010} & measured \\
16 & The other sensors (Table~\ref{tab:sensors}): smell has $\sim400$ receptor types, one per neuron, and a combinatorial code; taste works partly through ATP and serotonin; touch has fast- and slow-adapting sensors; warm and cold are separate & Mouse: one receptor recognizes many odors, one odor activates many receptors, and different odors activate different sets. Without ATP receptors the taste nerve does not respond; taste buds release serotonin. Single-fiber recordings from the skin of the human hand: four types by adaptation speed. The cold channel is distinct from the heat channels & \cite{Mombaerts2004,Malnic1999,Finger2005,Huang2005,JohanssonVallbo1979,McKemy2002} & measured \\
\bottomrule
\end{longtable}}
